\documentclass[10pt,a4paper]{amsart}
\usepackage[margin=19mm]{geometry}
\usepackage{amsmath,amssymb,mathtools}
\usepackage[scr=rsfs,cal=euler]{mathalfa}
\usepackage{xfrac}
\usepackage[unicode,hidelinks,bookmarks=false]{hyperref}
\newcommand{\ie}{i.e.\ }
\newcommand{\cf}{cf.\ }
\newcommand{\resp}{respectively\ }
\newcommand{\Subsection}{\subsection}
\providecommand{\nobreakdash}{\nobreak\hbox{-}}
\setlength{\emergencystretch}{2em}
\multlinegap0pt
\usepackage{amssymb}
\usepackage[shortlabels]{enumitem}
\usepackage{bbm}
\usepackage{multirow}
\usepackage{xcolor}
\usepackage{mathtools}
\usepackage{caption}
\usepackage[all]{xy}
\usepackage{tikz}
\newtheorem{thm}{Theorem}
\newtheorem{prop}{Proposition}
\newcommand{\F}{\mathbb{F}}
\DeclareMathOperator{\Gal}{Gal}
\newcommand{\Z}{\mathbb{Z}}
\title{Galois Groups of Apéry-like Series Modulo Primes}
\author{Xavier Caruso, Florian F\"urnsinn, Daniel Vargas-Montoya, Wadim Zudilin}
\date{Source: arXiv:2510.23298v1 (CC BY 4.0)}
\begin{document}
\maketitle

\noindent\emph{Source and licence.} Galois Groups of Apéry-like Series Modulo Primes. Version arXiv:2510.23298v1 (CC BY 4.0), distributed by arXiv under the Creative Commons Attribution 4.0 International licence, http://creativecommons.org/licenses/by/4.0/, as recorded in the arXiv licence metadata for this exact version and shown on the abstract page https://arxiv.org/abs/2510.23298v1. Full licence notice: \texttt{licenses/arxiv-2510.23298-CC-BY-4.0.txt}.\par\smallskip

\noindent\textit{Editorial scope.} Counted unit: the article's thm thm:apery (source0.tex lines 218--227). The article's complete original proof of it is delivered with it (source0.tex lines 461--477, one \texttt{proof} environment of 17 lines, which the article places away from the statement). No mathematics is written, repaired or re-derived: the statement and the proof are the article's own lines, in the article's own notation, and no retained line is substituted or re-flowed.  Retained uncounted as the source-local context the proof needs: lines 432--441.  Nothing was omitted from any retained span, so the package carries no omission record.  No retained line contains a citation, so the wrapper carries no bibliography and no bibliography entry.  No retained line contains a figure environment, an image-inclusion command or drawing code, so no image is delivered and the package declares no asset dependency.  Editorial formatting only, in addition: an AMS \texttt{amsart} preamble that restates the article's own declarations and macros used by the retained lines, namely the environments \texttt{thm}, \texttt{prop} on separate counters, together with the standard packages those lines need.  The retained environments print in this document as Theorem 1, Proposition 1 in order of appearance, whereas the article prints the same items with the numbering of its own full document.  The complete native source of the article as distributed in the arXiv e-print for this exact version is bundled unchanged as \texttt{sources/187870/source0.tex}.
\begin{prop} \label{prop:Galh}
    The Galois group of $\F_p(x,h^2)/\F_p(t)$ is 
    \[\Gal(\F_p(x,h^2)/\F_p(t)) = \begin{cases}
        \Z/(p{-}1)\Z\simeq \Z/e\Z\times \Z/2\Z & \text{if } p\equiv 3 \pmod 4\\
        \Z/(p{-}1)\Z & \text{if } p\equiv 13, 17 \pmod {24}\\
        \Z/e\Z\times \Z/2\Z& \text{if } p\equiv 1, 5 \pmod {24}.
    \end{cases}
    \]
    This corresponds to the cases where precisely one of the two numbers $u =\pm \frac 8 9$ corresponds to an involution $\sigma: h^2\mapsto u h^2 (1+x)^2$ in $\Gal(\F_p(x, h^2)/\F_p(t)$, none of them do, respectively both of them do.
\end{prop}

\begin{thm}~
\label{thm:apery}
\begin{itemize}
\item If $p \equiv 1, 5, 7, 11 \pmod{24}$, then
$\Gal\big(\F_p(t, f_\alpha)/\F_p(t)\big) = S$.

\item If $p \equiv 13, 17, 19, 23 \pmod{24}$, then
$\Gal\big(\F_p(t, f_\alpha)/\F_p(t)\big) = \F_p^\times$.
\end{itemize}
\end{thm}

\begin{proof}[Proof of Theorem~\ref{thm:apery}.]
    We distinguish cases according to the congruence class of $p$ modulo $24$:

    \begin{itemize}
        \item If $p\equiv 1 \pmod{24}$, the Galois group of $\F_p(x, h^2)/\F_p(t)$ is not cyclic, according to Proposition~\ref{prop:Galh}. Thus, $\F_p(t, f)$ is a proper subfield of $\F_p(x, h^2)$, necessarily of degree $e$ over $\F_p(t)$.
        \item If $p\equiv 5 \pmod{24}$ the Galois group of $\F_p(x, h^2)/\F_p(t)$ is not cyclic either, according to Proposition~\ref{prop:Galh}. We conclude again that $\F_p(f,t)$ has degree $e$ over $\F_p(t)$.
        \item If $p\equiv 7 \pmod{24}$, the Galois group of $\F_p(x, h^2)/\F_p(t)$ is cyclic, because $e$ is odd. Then, we need to check whether $f$ is in the unique subfield of $\F_p(x, h^2)$ of index 2. This is equivalent to checking whether it is invariant under the unique element  of the Galois group of order $2$. According to Proposition~\ref{prop:Galh}, this element $\sigma$ is uniquely given as the prolongation of $x\mapsto \frac{1-8x}{8+8x}$ with $\sigma(h^2)=u\cdot h^2\cdot (x+1)^2$, where the prefactor $u$ has to be $\frac89$. We have 
        \[\textstyle \sigma(f)= \sigma(h^2) (\sigma(x)+1) = \frac 8 9 h^2 (x+1)^2 (\sigma(x)+1) = h^2\cdot (x+1) = f.\]
        So indeed, $f$ is fixed by $\sigma$, and the extension has degree $e$.
        \item If $p\equiv 11 \pmod{24}$ one proceeds as for $7$, as $u=\frac 8 9$ in the definition of the involution $\sigma$.     
        \item If $p\equiv 13 \pmod{24}$, the group $\Gal(\F_p(x, h^2)/\F_p(t))$ is cyclic, and the unique element of order $2$ in the Galois group is an element of $S$, sending $h^2$ to $-h^2$. Thus $f$ is not fixed and $\F_p(t, f)\simeq \F_p(x, h^2)$, having degree $p{-}1$ over $\F_p(t)$.
        \item The case $p\equiv 17 \pmod{24}$ works as $p\equiv 13 \pmod{24}$, with the same conclusion.
        \item If $p\equiv 19 \pmod{24}$, we proceed as in the case $p\equiv 7 \pmod{24}$, with the exception that $u=-\frac 8 9$. Thus $\sigma(f)=-f$, and $f$ is not in the unique subfield of index $2$ of $\F_p(x, h^2)$. Thus $\F_p(t, f)$ has degree $p{-}1$.
        \item If $p\equiv 23 \pmod{24}$ we proceed as for $p\equiv 19 \pmod{24}$, as $u=-\frac 8 9$ again.
    \end{itemize}
This concludes the proof.
\end{proof}

\end{document}
